../cictikz/src/cictikz/data/tex/cictikz_preamble.tex

% Switched-capacitor amplifier, phase 1. C2 is shorted so V2 = 0, and C1 is
% not yet connected to the OTA, so the charge C1*V1 sits on C1 alone.
%
% Pairs with l5_scintro2, which is the same circuit with the two switches in
% the other state. They share tikz/sc_lib.tex, because the lecture's charge
% conservation argument works by comparing the two drawings.
%
% Polarity marks are the original's and they matter: + on C1's left plate,
% which is grounded, so its right plate goes to -V1 when the switch closes in
% phase 2, exactly as the prose says.

% Deviation from the original: the hand drawn sketch labels these
% "V_1 = l" and "V_2 =", with nothing after the equals sign. The
% dangling equals reads as a missing value rather than as a label, so
% the labels here are just the voltage names.

\begin{circuitikz}[american, thick, transform shape]

  ../cictikz/src/cictikz/data/tex/cictikz_lib.tex
  ../cictikz/src/cictikz/data/tex/sc_lib.tex

  \scIntroFrame

  % C1 not yet connected to the OTA
  \scSwOpen{2.6}{\scYin}
  % C2 shorted, so V2 = 0
  \scSwClosed{4.8}{\scYsh}

  \node[anchor=west] at (0.7,\scYin+1.05) {$V_{1}$};
  \node[anchor=west] at (4.1,\scYctwo-1.05) {$V_{2}$};

\end{circuitikz}
\end{document}
